\chapter{A Domain-Parametric Kernel for Safe Control Under Degraded Observation}
\label{ch:universality-completeness}

This chapter is the control-theory extension of the dissertation.  Its
purpose is not to add more runtime rows; it is to isolate the
domain-parametric kernel that the rest of the monograph instantiates.  The
central claim is deliberately narrow: under degraded observation, safe
control is governed by an observation-consistent state set, a
reliability-conditioned tightening law, a repair map, and a causal
certificate surface.  Battery remains the only full proof reference
surface.  Healthcare remains a defended instantiation, and autonomous
vehicles remain defended under a bounded contract within the active
three-domain program.

The point of the chapter is to make that claim precise enough that a reader
can separate three things cleanly: what the reference-row proof already establishes,
what the universal kernel adds, and where the current transfer boundary still
stops.  The implementation follows the same boundary.  The canonical
production kernel realizes the typed repair-and-certify contract, while the
formal arguments below isolate the control objects that must exist in any
domain instantiation.

In the paper-facing presentation, this chapter should be read as the final
piece of a promoted but scoped theory arc.  T11 is the main transfer theorem.
T9 and T10 are defended as observation-ambiguity lower bounds after the
theorem-promotion gates discharge their constants and domain assumptions for
the locked Battery, AV, and Healthcare artifacts.  T9 is the empty-safe-core
mandatory-release impossibility theorem, not a universal failure theorem for
controllers allowed to abstain or fail closed.  T10 is the two-state
total-variation lower bound, not a global minimax frontier.  \texttt{T\_minimax}
remains scoped to finite ambiguity classes; no unrestricted global minimax
optimality claim is made.

% ----------------------------------------------------------------
\section{Universal Problem Statement and Information Pattern}
\label{sec:universal-kernel-model}
% ----------------------------------------------------------------

We consider a partially observed, safety-constrained control system on a
discrete horizon $t = 0,1,\dots,T-1$:
\[
  x_{t+1} = f(x_t, a_t, \omega_t),
  \qquad
  y_t = h(x_t, \nu_t),
  \qquad
  w_t = \psi(y_{0:t}),
\]
where $x_t \in \mathcal{X} \subseteq \mathbb{R}^n$ is the true plant state,
$a_t \in \mathcal{A} \subseteq \mathbb{R}^m$ is the applied control,
$y_t$ is the telemetry channel output, $\omega_t$ is admissible plant
disturbance, $\nu_t$ is telemetry corruption, and $w_t \in [0,1]$ is the
reliability score extracted from the observation history.  The safe state set
is
\[
  \mathcal{S} = \{x \in \mathcal{X} : \phi(x) \ge 0\},
\]
and the violation count is
\[
  V_T = \sum_{t=0}^{T-1} \mathbf{1}[x_t \notin \mathcal{S}].
\]

The candidate controller sees only the observation history
$\mathcal{H}_t = \sigma(y_{0:t})$ and proposes
$a_t^{\mathrm{cand}} = \pi_t(\mathcal{H}_t)$.  The ORIUS kernel then builds
an observation-consistent state set $U_t$, derives a tightened safe-action
set $\mathcal{A}_t^{\mathrm{safe}}(U_t)$, repairs the candidate action into
$a_t^{\mathrm{safe}}$, and emits a certificate summarizing the causal data
used at step $t$.

\begin{definition}[Degraded-observation safe-control problem]
\label{def:degraded-observation-safe-control}
Given only the observation history $\mathcal{H}_t$, the reliability score
$w_t$, and a nominal candidate action $a_t^{\mathrm{cand}}$, construct a
measurable repair-and-certify map such that:
\begin{enumerate}
  \item the repaired action depends only on causal information available at
    time $t$;
  \item the repaired action preserves safety for all states consistent with
    the degraded observation, up to the stated miscoverage level $\alpha$;
  \item the resulting episode-level violation count admits an explicit bound
    in terms of $\alpha$, the reliability process, and the horizon.
\end{enumerate}
\end{definition}

For the control-theory reader, this chapter answers three questions beyond the
reference proof ladder:

\begin{table}[htbp]
\centering
\caption{Theorem questions answered by the monograph extension.}
\label{tab:universal-theorem-questions}
\begin{tabular}{p{1.5cm}p{4.7cm}p{5.1cm}}
\toprule
\textbf{Theorem} & \textbf{Control-theory question} & \textbf{Role in the monograph} \\
\midrule
T9 & When does observation-only mandatory release become impossible? & Empty-safe-core mandatory-release theorem \\
\addlinespace
T10 & How much risk is forced by indistinguishable boundary states? & Two-state TV lower-bound theorem \\
\addlinespace
T11 & Which structural obligations let the reference proof pattern transfer? & Template transfer theorem for new domains \\
\bottomrule
\end{tabular}
\end{table}

The dependency order is strict.  T1--T8 establish the reference proof
ladder.  T9 and T10 are assumption-qualified extensions under explicit
domain-discharge artifacts: T9 frames the structural risk of mandatory release
when the common safe core is empty, and T10 records a
boundary-indistinguishability lower bound for the stated binary testing model.
T11 then identifies the obligations under which
the same repair-and-bound structure transfers across domains.

\paragraph{Mathematical contract vs. domain discharge.}
This chapter separates theorem kernels from empirical validation.  The
\emph{mathematical contract} consists only of the displayed hypotheses,
Appendix~B assumptions, and typed transfer obligations: uncertainty-set
coverage, state-conditioned safe-set correctness, repair membership, fallback
admissibility, and domain-adapter validity.  The \emph{domain discharge}
artifacts instantiate constants and check mixing, boundary-mass, and witness
conditions for the promoted Battery, AV, and Healthcare runs.  In particular,
domain discharge artifacts are evidence, not an extra hidden assumption.
Consequently, the transfer claim is \emph{contract-universal, not unrestricted-global}: any physical-AI domain that implements the ORIUS runtime
contract may reuse the proof kernel, but each new domain must provide its own
discharge artifacts before numerical constants or deployment-facing claims are
promoted.

% ----------------------------------------------------------------
\section{Observation-Aware Control Is Structurally Necessary (T9)}
\label{sec:t9-universal-impossibility}
% ----------------------------------------------------------------

Theorem~\ref{thm:no-free-safety} shows, on the reference proof surface, that
quality-ignorant safety logic can fail even when observed-state evaluation
looks safe.  The broader universality claim is that this pathology is not a
single-domain accident; T9 records the corresponding scoped structural
statement for observation-only mandatory release on a nonempty ambiguity class
whose common safe core is empty.  The locked Battery, AV, and Healthcare
artifacts instantiate such ambiguity witnesses through the promotion evidence
files \texttt{T9\_battery.json}, \texttt{T9\_av.json}, and
\texttt{T9\_healthcare.json}.

\begin{theorem}[Impossibility of quality-ignorant mandatory release, T9]
\label{thm:universal-impossibility}
Let \(\Pi_{\mathrm{obs}}\) be total observation-only mandatory-release
controllers \(\pi:\mathcal O\to\mathcal A\), and let
\[
  B(o)=\{x\in\mathcal X:h(x)=o\},
  \qquad
  K(o)=\bigcap_{x\in B(o)}C(x)
\]
be the ambiguity class and common safe core for observation \(o\).  If
\(B(o)\ne\varnothing\) and \(K(o)=\varnothing\), then no controller
\(\pi\in\Pi_{\mathrm{obs}}\) can guarantee true-state safety for every
\(x\in B(o)\).
\end{theorem}

\begin{proof}
A total mandatory-release controller must choose one action \(a=\pi(o)\) for
the entire ambiguity class \(B(o)\). If that action were safe for every
\(x\in B(o)\), then
\[
  a\in\bigcap_{x\in B(o)}C(x)=K(o),
\]
contradicting \(K(o)=\varnothing\). Hence there exists
\(x^\dagger\in B(o)\) with \(a\notin C(x^\dagger)\). The release is therefore
not certifiably safe for the whole ambiguity class.
\end{proof}

\begin{remark}
T9 is deliberately narrower than the historical ``universal impossibility''
wording. It does not apply to policies allowed to abstain, expand the
uncertainty set, invoke fallback, or deny release. Those are precisely the
operations added by ORIUS. The theorem only attacks total observation-only
mandatory release on ambiguity classes with empty common safe core.
\end{remark}

T9 is a defended necessity extension for the universal kernel because it
identifies the exact structural condition under which observation-aware repair
or fail-closed behavior is necessary. Broader portability still requires
rerunning the ambiguity-class discharge builder on the new domain rather than
inheriting the Battery/AV/Healthcare witnesses.

% ----------------------------------------------------------------
\section{Boundary-Indistinguishability Lower Bound (T10)}
\label{sec:t10-lower-bound}
% ----------------------------------------------------------------

The core ORIUS upper bound is an envelope after a domain supplies a
per-step risk contract. T10 is the complementary lower-bound primitive: it
shows that if two boundary states induce nearly indistinguishable observations
while requiring disjoint safe actions, observation-only mandatory release must
incur nontrivial worst-state risk. The theorem is stated directly in total
variation distance. Any mapping from reliability \(w_t\) to a TV radius is a
domain-specific discharge step, not part of T10 itself.

\begin{theorem}[Boundary-indistinguishability lower bound, T10]
\label{thm:information-theoretic-lower-bound}
Let \(x_0,x_1\) be two latent boundary states with induced observation laws
\(P_0,P_1\). If
\[
  d_{\mathrm{TV}}(P_0,P_1)\le\epsilon
  \qquad\text{and}\qquad
  C(x_0)\cap C(x_1)=\varnothing,
\]
then every observation-only mandatory-release controller \(\pi\) satisfies
\[
  \sup_{i\in\{0,1\}}
  \Pr_{O\sim P_i}\!\left[\pi(O)\notin C(x_i)\right]
  \ge
  \frac{1-\epsilon}{2}.
\]
\end{theorem}

\begin{proof}
It suffices to prove the deterministic case; randomized controllers follow by
conditioning on the internal random seed and averaging. Define
\[
  S_i=\{o:\pi(o)\in C(x_i)\},\qquad i\in\{0,1\}.
\]
Because \(C(x_0)\cap C(x_1)=\varnothing\), the sets \(S_0\) and \(S_1\) are
disjoint. The average release success probability is
\[
  \frac12\{P_0(S_0)+P_1(S_1)\}.
\]
Since \(S_1\subseteq S_0^c\), the test that decides \(x_0\) on \(S_0\) and
\(x_1\) otherwise has success probability at least this release success. The
optimal testing success between \(P_0\) and \(P_1\) is
\((1+d_{\mathrm{TV}}(P_0,P_1))/2\le(1+\epsilon)/2\). Thus the average release
loss is at least \((1-\epsilon)/2\), and the larger of the two state-conditional
release risks is at least that average.
\end{proof}

\begin{remark}
T10 is an information-theoretic two-state lower bound. It does not claim a
globally sharp minimax constant, a horizon-summed frontier, or that the ORIUS
reliability score \(w_t\) is itself a TV distance. Domain artifacts may derive
a TV radius from a noise model and then instantiate T10, but that bridge is
separate from the theorem.
\end{remark}

The resulting reliability-risk discussion is therefore used conservatively. A
domain-specific model may derive a function \(\epsilon(w_t)\) and plug it into
the lower bound
\[
  \frac{1-\epsilon(w_t)}{2},
\]
but the manuscript does not identify this plug-in curve with a universal
frontier. The canonical runtime helper for the conservative upper envelope lives in
\path{src/orius/universal_theory/risk_bounds.py}.  The richer frontier analysis
used in this chapter is implemented in
\path{src/orius/universal/reliability_safety_frontier.py}.  Together these
helpers make the central design message operational: degraded observation
imposes a quantitative infrastructure requirement, not just a qualitative
warning.

% ----------------------------------------------------------------
\section{Typed Structural Transfer and Failure-Mode Converse (T11)}
\label{subsec:adapter-contract}
% ----------------------------------------------------------------

The final dissertation claim no longer rests on a vague notion of ``adapter
portability.''  The reusable objects are now explicit both in the writing and
in code.

\begin{definition}[Universal adapter contract]
\label{def:adapter-contract}
A domain instantiation is compatible with the degraded-observation transfer
claim if it supplies:
\begin{enumerate}
  \item an observation packet and parsed state;
  \item a reliability assessment $w_t \in [0,1]$;
  \item an observation-consistent state set $U_t$;
  \item a domain safety specification and tightened safe-action set
    $\mathcal{A}_t$;
  \item a repair operator returning a safe action or admissible fallback;
  \item a causal certificate recording the repaired action, uncertainty state,
    reliability state, and assumption tags.
\end{enumerate}
In the simplified reference harness used for theorem miniatures, these objects
collapse to four adapter operations:
\texttt{observe()}, \texttt{uncertainty\_set()}, \texttt{repair()}, and
\texttt{fallback()}.
\end{definition}

The canonical typed implementation is provided by
\path{src/orius/universal_theory/contracts.py} and
\path{src/orius/universal_theory/kernel.py}.  The smaller reference harness in
\path{src/orius/universal/contract.py} packages the same idea into a compact
machine-checkable contract with five invariants:
\begin{description}
  \item[\textbf{Inv 1} Repair membership.] The repaired action lies in the
    tightened safe-action set.
  \item[\textbf{Inv 2} Monotone tightening.] Lower reliability never enlarges
    the admissible set.
  \item[\textbf{Inv 3} Reliability boundedness.] Reliability remains in
    $[0,1]$.
  \item[\textbf{Inv 4} Calibration consistency.] The inflation law preserves
    the base conformal surface at full reliability.
  \item[\textbf{Inv 5} Failure completeness.] Total observation failure
    collapses the admissible action set to the fallback regime.
\end{description}

\begin{theorem}[Typed structural transfer theorem, T11]
\label{thm:universality-completeness}
Let $\mathcal{D}$ be a domain instantiation for the degraded-observation
safe-control problem in Definition~\ref{def:degraded-observation-safe-control}.
Assume that, for every step $t$:
\begin{enumerate}
  \item the constructed observation-consistent state set satisfies
    $\Prob[x_t \in U_t \mid \mathcal{H}_t] \ge 1-\alpha$;
  \item the tightened safe-action set is sound for that uncertainty set, in
    the sense that every action in $\mathcal{A}_t^{\mathrm{safe}}(U_t)$ keeps
    the next state inside $\mathcal{S}$ for all $x \in U_t$ and all
    admissible disturbance;
  \item the repair map returns an element of
    $\mathcal{A}_t^{\mathrm{safe}}(U_t)$ whenever that set is non-empty;
  \item when the tightened safe-action set is empty, the fallback action is
    admissible and preserves safety from any current safe state.
\end{enumerate}
Then the repaired action produced by the ORIUS kernel satisfies the same
one-step safety-preservation statement as the reference theorem:
\[
  \Prob[x_{t+1}\in\mathcal{S}\mid \mathcal{H}_t] \ge 1-\alpha.
\]
\noindent In this theorem the four listed obligations are the entire forward
transfer surface.  No broader episode-level claim is inherited unless the new
domain also supplies an explicit predictable per-step risk budget as in
Corollary~\ref{cor:t11-aggregation}.  The converse remains a separate
failure-mode proposition rather than part of T11 itself.
\end{theorem}

\begin{proof}
Fix a step $t$ and condition on the observation history $\mathcal{H}_t$.
On the event $E_t = \{x_t \in U_t\}$, Assumption~(2) implies that every action
in $\mathcal{A}_t^{\mathrm{safe}}(U_t)$ preserves membership in the safe set
at the next step.  If the tightened safe-action set is non-empty, then
Assumption~(3) ensures that the repaired action returned by the kernel lies in
that set, hence preserves safety on $E_t$.  If the tightened safe-action set
is empty, then Assumption~(4) ensures that the fallback action preserves
safety from any current safe state.  Therefore,
\[
  \Prob[x_{t+1}\in\mathcal{S} \mid \mathcal{H}_t]
  \;\ge\;
  \Prob[E_t \mid \mathcal{H}_t]
  \;\ge\;
  1-\alpha,
\]
which is the one-step safety-preservation statement.
\end{proof}

\begin{corollary}[T10\_T11\_ObservationAmbiguitySandwich: Covered Observation-Ambiguity Safety Sandwich]
\label{cor:observation-ambiguity-sandwich}
Let $h:\mathcal{X}\to\mathcal{O}$ be the observation map.  For an observation
$o$, define the ambiguity class
\[
  B(o)=\{x\in\mathcal{X}:h(x)=o\}
\]
and the common safe core
\[
  K(o)=\bigcap_{x\in B(o)}C(x).
\]
For every mandatory-release observation-only controller
$\pi:\mathcal{O}\to\mathcal{A}$,
\[
  \operatorname{TSVR}(\pi)
  \;\ge\;
  \mathbb{E}_{O}\!\left[
    \min_{a\in\mathcal{A}}
    \Prob(a\notin C(X^\star)\mid O)
  \right].
\]
In particular, if $K(o)=\emptyset$ on a positive-probability ambiguity class,
no observation-only mandatory-release controller can guarantee zero
true-state violation on that class.  Conversely, if ORIUS constructs
$\mathcal{X}_t$ with
$\Prob(X_t^\star\notin\mathcal{X}_t)\le\alpha$ and releases only actions
\[
  u_t\in\bigcap_{x\in\mathcal{X}_t}C(x),
\]
then
\[
  \operatorname{TSVR}_{\mathrm{ORIUS}}\le\alpha,
\]
and under exact coverage $X_t^\star\in\mathcal{X}_t$ the one-step
true-state violation probability is zero.
\end{corollary}

\begin{proof}
Condition on an observation $O=o$.  Any observation-only mandatory-release
controller must choose a single action $a=\pi(o)$ for all states in $B(o)$.
Its conditional violation probability is
$\Prob(a\notin C(X^\star)\mid O=o)$, and the best possible observation-only
choice is therefore the displayed minimum over $a\in\mathcal{A}$.  Taking
expectation over $O$ gives the lower bound.  If $K(o)=\emptyset$, every
candidate action is unsafe for at least one state in $B(o)$, so the
conditional lower bound is positive whenever that unsafe state has positive
conditional mass.

For the ORIUS upper bound, split the violation event according to whether
$X_t^\star\in\mathcal{X}_t$.  On the coverage event, the released action lies
in $C(x)$ for every $x\in\mathcal{X}_t$ and hence in $C(X_t^\star)$, so the
violation probability on that event is zero.  The only remaining probability
mass is the coverage-miss event, bounded by $\alpha$.
\end{proof}

\paragraph{Promoted domain runtime lemmas.}
The AV and healthcare rows do not promote a new universal theorem.  They are
supporting runtime lemmas under T11: AV is bounded to the emitted
brake-hold postcondition, and healthcare is bounded to the emitted fail-safe
alert-release postcondition.  In both cases, missing T11 status, failed T11
obligations, invalid certificates, or false true-state postconditions fail the
runtime witness closed.

\begin{corollary}[Episode aggregation under explicit per-step budgets]
\label{cor:t11-aggregation}
If, in addition, the instantiated domain supplies a predictable per-step risk
budget sequence $r_t$ satisfying
\[
  \Prob[x_{t+1}\notin\mathcal{S}\mid \mathcal{H}_t] \le r_t
\]
for all $t$, then
\[
  \mathbb{E}[V_T] \le \sum_{t=1}^T r_t.
\]
When a particular domain instantiation justifies the reference-row envelope
$r_t \le \alpha(1-w_t)$, the episode bound
$\mathbb{E}[V_T] \le \alpha(1-\bar{w})T$ follows immediately.
\end{corollary}

\begin{proof}
Let $Z_t=\mathbf{1}\{x_{t+1}\notin\mathcal{S}\}$.  The assumed bound implies
\[
  \mathbb{E}[Z_t \mid \mathcal{H}_t] \le r_t.
\]
Taking expectations and summing over $t$ gives
\[
  \mathbb{E}[V_T]
  =
  \sum_{t=1}^T \mathbb{E}[Z_t]
  \le
  \sum_{t=1}^T r_t.
\]
The reference-row envelope is the special case $r_t \le \alpha(1-w_t)$.
\end{proof}

\begin{proposition}[Failure of any transfer obligation breaks the reference proof pattern]
\label{prop:t11-counterexamples}
If any one of the four T11 obligations fails, then there exists a domain
instantiation for which the reference transfer argument breaks.
\end{proposition}

\begin{proof}
If obligation~(1) fails, the latent state may lie outside $U_t$ with
probability larger than $\alpha$, so the repair step can be safe only for the
wrong uncertainty set and the one-step safety argument no longer applies.
If obligation~(2) fails, an action may belong to the purported tightened set
while still allowing a successor outside $\mathcal{S}$ for some state in
$U_t$, so soundness is lost directly.  If obligation~(3) fails, the repair map
can return an action outside the tightened safe set even when a safe repaired
action exists, which again breaks the one-step argument.  If obligation~(4)
fails, the kernel can reach a state where the tightened safe-action set is
empty and no admissible fallback exists, so the transfer proof has no safe
release path.  In each case the failure is structural rather than rhetorical:
the missing obligation removes one of the exact proof steps used in the T11
forward direction.
\end{proof}

T11 is therefore a typed structural transfer theorem paired with an explicit
failure-mode converse, not a blanket declaration that every evaluated row
already satisfies those obligations.  The energy-management reference row provides the full reference
proof.  Other domains enter the monograph only when their instantiation
evidence is strong enough to support the corresponding one-step obligations,
and the stronger episode-level claim appears only when a domain-specific
per-step risk budget is justified explicitly.

\subsection{Machine-Checkable Contract Surfaces}

Two levels of executable checking now exist in the repository.

\begin{enumerate}
  \item \textbf{Production structural contract.}
    \path{src/orius/universal_theory/contracts.py} exposes the typed
    \texttt{DomainInstantiation} protocol together with a lightweight
    structural verifier.  The focused regression test
    \path{tests/test_universal_theory.py} checks canonical battery,
    healthcare and vehicle adapters against that surface and then
    exercises the new universal step kernel end-to-end.
  \item \textbf{Reference invariant harness.}
    \path{src/orius/universal/contract.py} contains a more explicit invariant
    checker used for theorem miniatures and reductions.  The focused tests
    \path{tests/test_universal_contract.py} and
    \path{tests/test_unification.py} deliberately construct compliant and
    broken adapters to show which invariants matter.
\end{enumerate}

The simplified usage pattern is:

\begin{verbatim}
from orius.universal.contract import ContractVerifier

verifier = ContractVerifier(adapter, alpha=0.05)
verifier.check(z_t=np.array([0.0]), q_t=0.1)
\end{verbatim}

That harness is intentionally modest.  It does not replace the defended
domain-specific artifacts.  It makes the transfer theorem executable: the
reader can see the exact invariants, violate them, and observe the resulting
contract failure.

% ----------------------------------------------------------------
\section{Unification: ORIUS Subsumes Prior Frameworks}
\label{sec:unification}
% ----------------------------------------------------------------

\begin{theorem}[Unification]
\label{thm:unification}
The following prior safe-control families can be represented as special cases
of the ORIUS degraded-observation kernel:
\begin{enumerate}
  \item \textbf{Control Barrier Functions (CBF).} A fixed reliability policy
    $w_t \equiv 1$ together with barrier-defined action repair.
  \item \textbf{Robust MPC.} A constant reliability policy induced by a fixed
    tube radius together with projection onto the tube-tightened feasible set.
  \item \textbf{Safe RL.} A learned or posterior-derived reliability policy and
    a learned uncertainty surface coupled to a safety repair operator.
\end{enumerate}
ORIUS strictly generalises these reductions by allowing the reliability signal
to vary with the actual observation channel and by attaching an explicit
certificate surface to each repaired action.
\end{theorem}

\begin{proof}[Proof sketch]
The repository contains executable reference reductions in
\path{src/orius/universal/unification.py}.  \texttt{CBFAsORIUS} fixes
$w_t \equiv 1$ and uses barrier-style projection.  \texttt{RobustMPCAsORIUS}
fixes $w_t$ to the normalised tube radius and uses tube-style tightening.
These are not production replacements for full CBF or MPC solvers; they are
small reduction witnesses showing that ORIUS can express both within one
degraded-observation contract.  Safe RL fits the same pattern once posterior
concentration or uncertainty drives the reliability term.
\end{proof}

The unification is summarised in Table~\ref{tab:unification-comparison}.

\begin{table}[htbp]
\centering
\caption{Existing safe-control families as ORIUS reductions.}
\label{tab:unification-comparison}
\begin{tabular}{p{2.8cm}p{2.8cm}p{3.5cm}p{3.8cm}}
\toprule
\textbf{Framework} & \textbf{Reliability policy} & \textbf{Uncertainty / repair surface} & \textbf{ORIUS generalises by} \\
\midrule
CBF & $w_t \equiv 1.0$ & Barrier-defined safe set and projection & Adapting safety margins when telemetry degrades \\
\addlinespace
Robust MPC & $w_t \equiv c$ & Fixed tube tightening and projection & Using time-varying reliability instead of one constant tube level \\
\addlinespace
Safe RL & $w_t$ from posterior concentration & Learned uncertainty plus safety repair & Replacing model-specific confidence with an explicit degraded-observation contract \\
\addlinespace
\textbf{ORIUS / DC3S} & dynamic $w_t$ from OQE & Observation-consistent state set, repair, fallback, and certificate & \textit{The reusable degraded-observation kernel} \\
\bottomrule
\end{tabular}
\end{table}

% ----------------------------------------------------------------
\section{Constraint-Class Mismatch: When 1-D Action Repair Is Insufficient}
\label{sec:constraint-class-mismatch}
% ----------------------------------------------------------------

The DC3S shield projects each candidate action onto a tightened interval
derived from the conformal uncertainty set.  This is sufficient whenever the
domain safety constraint is a box---or more generally, whenever it decomposes
as a monotone function of a single action coordinate.  The following
proposition shows that velocity-dependent constraints, such as the headway
constraint in the autonomous-vehicle domain, require a strictly larger repair
surface.

\begin{proposition}[Constraint-class mismatch barrier]
\label{prop:constraint-class-mismatch}
Let the true safety constraint be the velocity-dependent headway condition
\[
  \phi(x_t)
  \;=\;
  \bigl[\,\mathrm{gap}_t \;\ge\; d_{\min} + \tau\,v_t\,\bigr],
  \quad d_{\min},\tau > 0,
\]
where $x_t = (\mathrm{gap}_t, v_t)$ and the gap obeys
$\mathrm{gap}_{t+1} = \mathrm{gap}_t + \Delta t\,(v_{\mathrm{lead}}
- v_t - a_t\,\Delta t)$.  Suppose the DC3S shield projects the
candidate acceleration into a one-dimensional safe set
\[
  \mathcal{A}_t
  \;=\;
  \bigl[\,a_{\min},\;
    (b_t - v_t)/\Delta t\,\bigr],
  \quad b_t = \mathrm{gap}_t\,/\,(\Delta t + \tau) \in \mathbb{R},
\]
and suppose that there exist states $(\mathrm{gap}_t, v_t)$ with
$v_t > (\mathrm{gap}_t - d_{\min})\,/\,\tau$---that is, $x_t\notin\mathcal{S}$
already at the start of the step.
Then:
\begin{enumerate}
  \item No action $a_t \in [a_{\min}, a_{\max}]$ moves $x_{t+1}$ into
    $\mathcal{S}$ in one step, because the constraint violation is
    determined by \emph{current} $(v_t, \mathrm{gap}_t)$, not by the
    next-step speed the shield controls.
  \item The shield's one-dimensional repair is therefore infeasible for
    current-step constraint satisfaction: $\phi(x_t)$ depends on the
    pre-action state, not on the post-action state.
  \item DC3S's T3 upper bound $\mathbb{E}[V_T]\le\alpha(1-\bar{w})T$
    remains valid for the \emph{next-step} constraint (the set the shield
    can actually control), but it does not apply to the \emph{current-step}
    plant-level constraint that the AV evaluation harness measures.
\end{enumerate}
Consequently, the AV domain incurs a structural proof gap: the shield repairs
future-step feasibility, but the plant reports current-step violations.
\end{proposition}

\begin{proof}
If $v_t > (\mathrm{gap}_t - d_{\min})\,/\,\tau$, then
$\mathrm{gap}_t < d_{\min} + \tau v_t$, so $\phi(x_t) = \mathrm{False}$
regardless of $a_t$.  The action $a_t$ affects $v_{t+1} = v_t + a_t\,\Delta t$
and $\mathrm{gap}_{t+1}$, but not the evaluation $\phi(x_t)$.
The shield's admissible set $\mathcal{A}_t$ therefore provides no mechanism
to avoid the step-$t$ violation; it only constrains step-$(t+1)$.
\end{proof}

\begin{corollary}[AV promotion routes from the original headway mismatch]
\label{cor:av-promotion}
Discharging the original autonomous-vehicle headway mismatch requires at
least one of the following:
\begin{enumerate}[label=(\alph*)]
  \item \emph{Augmented state repair}: the shield solves a two-dimensional
    quadratic program over $(a, v)$ that explicitly enforces
    $\mathrm{gap}_t \ge d_{\min} + \tau v_t$ as a current-state feasibility
    constraint, activating a full-brake fallback when the state is already
    outside~$\mathcal{S}$.
  \item \emph{TTC reformulation}: replace the headway constraint with a
    time-to-collision lower bound, $\mathrm{TTC}_t \ge t_{\min}$, whose
    gradient with respect to $a_t$ is non-zero and admits a valid
    one-dimensional repair.
  \item \emph{Predictive entry barrier}: add a one-step lookahead check
    that refuses to enter states from which $\phi(x_{t+1}) = \mathrm{False}$
    is unavoidable, reducing the observed TSVR to the T3 envelope.
\end{enumerate}
Without one of these routes, the AV residual is a \emph{structural} residual,
not a conformal-coverage residual.
\end{corollary}

\paragraph{Current closure route.}
This dissertation adopts the TTC reformulation together with the predictive
entry barrier as the canonical AV promotion path.  The two-dimensional
$(a,v)$ repair remains a valid future route, but it is not the route claimed
here.  Under the bounded TTC contract, the closure matrix now records the AV
row as proof-validated: locked replay reduces TSVR to zero on the canonical
validation harness, the conservative entry-barrier soundness cases clear, and
the CertOS lifecycle executes on the same adapter surface.  This promotion is
deliberately narrow.  It does not promote every richer headway, lane-change,
or multi-agent vehicle setting to deeper proof depth within the same bounded
three-domain program.

% ----------------------------------------------------------------
\section{Chapter Summary}
\label{sec:ch37-summary}
% ----------------------------------------------------------------

This chapter closes the monograph at the correct control-theory level.

\begin{enumerate}
  \item \textbf{T9} is a supporting extension that frames observation-aware
    safety logic as structurally necessary under persistent degradation once
    the domain mixing and witness constants are discharged.
  \item \textbf{T10} is a supporting extension that turns mean reliability
    into a quantitative infrastructure-to-risk frontier under explicit
    binary-testing and boundary-law hypotheses.
  \item \textbf{T11} packages universality as a typed kernel and adapter
    template, not as a claim that every domain already has battery-level proof
    closure.
  \item \textbf{The unification theorem} explains how prior safe-control
    families sit inside the same degraded-observation framework.
\end{enumerate}

The central theorem contribution is now precise: the safe-control object under
degraded telemetry is not the nominal state alone but a reliability-indexed
observation-consistent state set together with a repair map and a causal
certificate surface.  Battery supplies the full proof witness.  T9 and T10,
as assumption-qualified supporting extensions, record necessity and lower-bound
directions after the locked domains discharge the explicit constants.  T11
shows how the same kernel transfers when a new domain discharges the explicit
structural obligations.  This is a defensible control-theory claim; it is not
an assertion that all evaluated domains already enjoy equal proof closure.
